
\section{Experimental Results}
\begin{table*}[t]\centering
\vspace{-10pt}
\caption{Comparison against baselines on TL;DR, with target agreement level $1 - \alpha = 0.9$. The results are averaged across 1000 runs with random data split. Guarantee Success Rate is defined as the ratio of successful runs that achieve empirical human agreement larger than or equal to $1 - \alpha$. \textbf{Cascaded Selective Evaluation is the only method that achieves high guarantee success rate, while maintaining high coverage.}}
\vspace{-5pt}
\resizebox{.96\textwidth}{!}{
    \begin{tabular}{lcccccc}\toprule
         \multirow{2}{*}[-0.3em]{\textbf{Method}} & \multicolumn{3}{c}{\textbf{Evaluator Composition (\%)}} & \multirow{2}{*}[-0.3em]{\textbf{Coverage (\%)}} & \multirow{2}{*}[-0.3em]{\shortstack{\textbf{Guarantee Success}\\\textbf{Rate (\%)}}}\\\cmidrule(lr){2-4}
         & \textit{Mistral-7B} & \textit{GPT-3.5-turbo} & \textit{GPT-4-turbo} & & \\
         \midrule
         No Selection & 0.0 & 0.0 & 100.0 & 100.0 & 0.0 \\
         Heuristic Selection & 0.0 & 0.0 & 100.0 & 89.6 & 42.0 \\
         Cascaded Heuristic Selection & 59.6 & 15.0 & 25.5 & 64.6 & 0.0 \\
         \midrule
         \multirow{3}{*}{Point-Estimate Calibration} & 100.0 & 0.0 & 0.0 & 5.6 & 54.7 \\
                                                     & 0.0 & 100.0 & 0.0 & 9.4 & 79.0 \\
                                                     & 0.0 & 0.0 & 100.0 & 57.7 & 47.5 \\
         \midrule
         \textbf{Cascaded Selective Evaluation} & \textbf{28.3} & \textbf{28.2} & \textbf{43.5} & \textbf{55.7} & \textbf{90.8} \\
         \bottomrule
    \end{tabular}
}
\vspace{-5pt}
\label{tab:tldr_results}
\end{table*}
\begin{figure*}[t]
    \centering
    \includegraphics[width=.96\textwidth]{figures/tldr_results.pdf}
    \vspace{-5pt}
    \caption{TL;DR results. \textbf{Cascaded Selective Evaluation guarantees human agreement far beyond a level achievable by GPT-4 without abstention (Left), while employing substantially weaker judge models (Right).} Solid blue line denotes average human agreement over 1000 runs on the dataset, and the light blue region denotes the min / max agreement within the 1000 runs.}
    \vspace{-10pt}
    \label{fig:tldr_results}
\end{figure*}
\subsection{Evaluating Generated Summaries}\label{sec:evaluating_summaries_on_TLDR}
\paragraph{Experimental Setup.} We first test our approach for evaluating summaries on TL;DR dataset \citep{summarize-from-human-feedback}. We use a cascade of \textit{Mistral-7B-instruct-v0.2} \citep{mistral}, \textit{GPT-3.5-turbo} and \textit{GPT-4-turbo} \citep{gpt-4} as judges. Observing that the dataset provides multiple human annotations per input, we use \textit{Simulated Annotators (Ind.)} with $K = N = 5$. We fix the size of calibration set $|D_\textit{cal}| = 500$ and $\delta = 0.1$, and run the experiments for 1000 random splits of calibration and test set. For baselines, we consider: (1) Heuristic Selection: using GPT-4 as a judge and setting $\lambda = 1 - \alpha$, assuming perfect calibration; (2) Cascaded Heuristic Selection: a variant of Heuristic Selection using the same cascades of judge models as ours; (3) Point-Estimate Calibration: setting $\lambda$ as the smallest value that satisfies $\widehat{R}(\lambda) \le \alpha$ in $D_\textit{cal}$, without hypothesis testing.

In Figure \ref{fig:tldr_results}, we show that human agreement guarantee is satisfied with our approach across all levels of target human agreement, far beyond what GPT-4 can achieve without abstention. Notably, unlike prior works \citep{conformal-alignment, conformal-factuality-guarantee} that only controls the risk in expectation over calibration sets (solid blue line), our method guarantees with high probability that each individual run would satisfy the target agreement level (light blue region). Moreover, as shown in right plot, the high agreements can be achieved while the majority of evaluation are done with substantially smaller LLMs than GPT-4. For example, our method can outperform GPT-4 with 80\% human agreement, while 75\% of the evaluations are done by Mistral-7B or GPT-3.5. 

We also compare against selective baselines in Table \ref{tab:tldr_results}, in terms of their coverage and guarantee success rate. All baselines fail to provide meaningful guarantee success rate without significantly sacrificing the coverage. This includes Point-Estimate Calibration, which makes use of the test statistics in calibration data. On the contrary, \method achieves over 90\% success rate---just as expected by setting $\delta = 0.1$---attesting to its reliability.

\begin{table*}[t]\centering
\vspace{-15pt}
\caption{Comparison to baselines on ChatArena, with target agreement level $1 - \alpha = 0.85$. The results are averaged across 1000 runs with random data split. Consistent with results on TL;DR, \textbf{our method successfully guarantees target agreement level while maintaining high coverage.}}
\vspace{-5pt}
\resizebox{.96\textwidth}{!}{
    \begin{tabular}{lcccccc}\toprule
         \multirow{2}{*}[-0.3em]{\textbf{Method}} & \multicolumn{3}{c}{\textbf{Evaluator Composition (\%)}} & \multirow{2}{*}[-0.3em]{\textbf{Coverage (\%)}} & \multirow{2}{*}[-0.3em]{\shortstack{\textbf{Guarantee Success}\\\textbf{Rate (\%)}}}\\\cmidrule(lr){2-4}
         & \textit{Mistral-7B} & \textit{GPT-3.5-turbo} & \textit{GPT-4-turbo} & & \\
         \midrule
         No Selection & 0.0 & 0.0 & 100.0 & 100.0 & 0.0 \\
         Heuristic Selection & 0.0 & 0.0 & 100.0 & 95.2 & 0.1 \\
         Cascaded Heuristic Selection & 57.1 & 15.2 & 27.7 & 79.7 & 0.3 \\
         \midrule
         \multirow{3}{*}{Point-Estimate Calibration} & 100.0 & 0.0 & 0.0 & 0.0 & 0.0 \\
                                                     & 0.0 & 100.0 & 0.0 & 40.5 & 57.2 \\
                                                     & 0.0 & 0.0 & 100.0 & 60.9 & 54.4 \\
         \midrule
         \textbf{Cascaded Selective Evaluation} & \textbf{23.7} & \textbf{58.8} & \textbf{17.5} & \textbf{63.2} & \textbf{91.0} \\
         \bottomrule
    \end{tabular}
}
\vspace{-5pt}
\label{tab:chatarena_results}
\end{table*}
\begin{figure*}[t]
    \centering
    \includegraphics[width=.96\textwidth]{figures/chatarena_results.pdf}
    \vspace{-5pt}
    \caption{ChatArena results. \textbf{Our approach guarantees target human agreement level (Left) while majority of evaluations are done with weaker judge models, Mistral-7B and GPT-3.5 (Right).}}
    \vspace{-15pt}
    \label{fig:chatarena_results}
\end{figure*}

\subsection{Evaluating LLM-based Chat Assistants} \label{sec:evaluating_chat_assistants}
\paragraph{Experimental Setup.} Next, we test our approach for evaluating general-purpose LLM assistants on two datasets: Chat(bot) Arena \citep{mt-bench} with real-world user-assistant interaction\footnote{We use an evaluation set with 5.2k instances sampled by \cite{preference-dissection}.} and Auto-J \citep{auto-j}, a curated benchmark for meta-evaluation of LLM judges. We employ the same cascades of models as in \S \ref{sec:evaluating_summaries_on_TLDR}, but this time we use \textit{Simulated Annotators (Maj.)}, as both datasets only provide one human annotation per input $x$. We set $K = N = 5$ and $|D_\textit{cal}| = 500$ for ChatArena, and $K = N = 3$, $|D_\textit{cal}| = 392$ for Auto-J, considering the small size of the benchmark.

The results are shown in Figure \ref{fig:chatarena_results} and Table \ref{tab:chatarena_results} for ChatArena, and in \S \ref{app:additional_results_on_chat_assistant_evaluation} for Auto-J. Again, we confirm that human agreement guarantee can be achieved across all levels of $\alpha$. On ChatArena, unlike Point-Estimate Calibration with GPT-4 whose guarantee success rate is below 60\%, our method achieves 91\% guarantee success rate while only using GPT-4 for 17.5\% of the evaluations. The performance is particularly pronounced in Auto-J where GPT-4 without abstention could only achieve 63.2\% agreement with humans (Figure \ref{fig:auto-j_results}), potentially due to the fact that the dataset introduces additional \textit{tie} label unlike the other two datasets. In stark contrast, \method guarantees up to 80\% human agreement with high probability.

Next, we conduct further ablations and analyses to better understand the working of our method.
\vspace{-5pt}

\subsection{Understanding Abstention Policy}
\begin{wrapfigure}{r}{.48\textwidth}
    \centering
    \renewcommand{\arraystretch}{1.0}
    \small
    \vspace{-10pt}
    \captionof{table}{Comparison between abstained vs. evaluated samples. \textbf{Our abstention policy aligns with how humans agree with each other (IAA)}, exhibiting no significant reliance on shallow heuristics (length ratio, token overlap).}
    \resizebox{.48\textwidth}{!}{
    \begin{tabular}{ lcc }\toprule
        \textbf{Dimension} & \textbf{Abstained Samples} & \textbf{Evaluated Samples} \\
        \midrule
        Human IAA & 0.815 (0.031) & 0.902 (0.025) \\
        Length Ratio &  0.242 (0.014) & 0.245 (0.025) \\
        Token Overlap & 0.623 (0.049) & 0.592 (0.054) \\
        \bottomrule
    \end{tabular}
    }
    \vspace{-5pt}
    \label{tab:analyzing_abstention_policy}
\end{wrapfigure}
One potential concern with selective evaluation is that its abstention policy might not align with the human-perceived subjectivity of each instance and instead rely on shallow heuristics, \eg choosing to abstain when the pair of generations have large token overlap. To address this concern, we analyze whether there exists a significant difference in human-perceived subjectivity between model-abstained samples and evaluated samples. We first collect 3-5 human annotations per each instance in ChatArena, and measure the inter-annotator agreement\footnote{As the label space is binary for preference evaluation, we define inter-annotator agreement simply as the density of the majority preference label assigned by human annotators.} (IAA) as a proxy of human-perceived subjectivity. Then, we compare the difference in IAA between abstained and evaluated samples when the target agreement level is set to 0.9. We also measure the difference in terms of shallow features, specifically the pairwise length ratio and the token-overlap (ROUGE-L) within each instance. For further details, see \S \ref{app:human_evaluation_details}.

\begin{table*}[t]\centering
\vspace{-15pt}
\caption{Impact of number of simulated annotators $N$ on ChatArena, with $1 - \alpha = 0.85$. Larger number of simulations generally leads to better coverage, while human agreement is guaranteed even with a small $N$. \textbf{For all values of $N$, Cascaded Selective Evaluation guarantees high agreement with humans while reducing the API cost by 40\% compared to GPT-4 without abstention.}}
\vspace{-5pt}
\resizebox{.93\textwidth}{!}{
    \begin{tabular}{lcccc}\toprule
         \multirow{2}{*}{\textbf{Method}} & \multirow{2}{*}{\shortstack{\textbf{Empirical Human}\\\textbf{Agreement (\%)}}} & \multirow{2}{*}{\textbf{Coverage (\%)}} & \multirow{2}{*}{\shortstack{\textbf{Guarantee Success}\\\textbf{Rate (\%)}}} & \multirow{2}{*}{\shortstack{\textbf{Relative }\\\textbf{API Cost}}} \\
         & & & & \\
         \midrule
         % \rowcolor{gray!25}\multicolumn{5}{c}{$\mathbf{N = 0}$ (Zero-shot Predictive Probability)}\\
         GPT-4 ($N = 1$) & 77.8 & 100.0 & 0.0 & 1.000 \\
         Cascaded Selective Evaluation ($N = 1$) & 85.2 & 60.9 & 90.3 & 0.655 \\
         \midrule
         % \rowcolor{gray!25}\multicolumn{5}{c}{$\mathbf{N = 2}$}\\
         GPT-4 ($N = 2$) & 78.2 & 100.0 & 0.0 & 2.000 \\
         Cascaded Selective Evaluation ($N = 2$) & 85.7 & 61.5 & 90.8 & 1.288 \\
         \midrule
         % \rowcolor{gray!25}\multicolumn{5}{c}{$\mathbf{N = 3}$}\\
         GPT-4 ($N = 3$) & 78.1 & 100.0 & 0.0 & 3.000 \\
         Cascaded Selective Evaluation ($N = 3$) & 85.5 & 62.1 & 90.3 & 1.920 \\
         \midrule
         % \rowcolor{gray!25}\multicolumn{5}{c}{$\mathbf{N = 5}$}\\
         GPT-4 ($N = 5$) & 78.5 & 100.0 & 0.0 & 5.000 \\
         Cascaded Selective Evaluation ($N = 5$) & 85.8 & 63.2 & 91.0 & 2.849 \\
         \bottomrule
    \end{tabular}
}
\vspace{-10pt}
\label{tab:impact_of_number_of_simulated_annotators}
\end{table*}

% 

% \begin{table*}[t]\centering
% \vspace{-15pt}
% \caption{Impact of number of simulated annotators $N$ on ChatArena, with $1 - \alpha = 0.85$. Larger number of simulations generally leads to better coverage, while human agreement is guaranteed even with a small $N$. \textbf{At all levels of $N$, Cascaded Selective Evaluation provides higher agreement with humans than GPT-4 while reducing the API cost by 40\%.}}
% \vspace{-5pt}
% \resizebox{.93\textwidth}{!}{
%     \begin{tabular}{clcccc}\toprule
%          \multirow{2}{*}{$\mathbf{N}$} & \multirow{2}{*}{\textbf{Method}} & \multirow{2}{*}{\shortstack{\textbf{Empirical Human}\\\textbf{Agreement (\%)}}} & \multirow{2}{*}{\textbf{Coverage (\%)}} & \multirow{2}{*}{\shortstack{\textbf{Guarantee Success}\\\textbf{Rate (\%)}}} & \multirow{2}{*}{\shortstack{\textbf{Relative }\\\textbf{API Cost}}} \\
%          & & & & \\
%          \midrule
%          % \rowcolor{gray!25}\multicolumn{5}{c}{$\mathbf{N = 0}$ (Zero-shot Predictive Probability)}\\
%          \multirow{2}{*}{1} & GPT-4 & 77.8 & 100.0 & 0.0 & 1.000 \\
%          & Cascaded Selective Evaluation & 85.0 & 63.8 & 90.3 & 0.655 \\
%          \midrule
%          % \rowcolor{gray!25}\multicolumn{5}{c}{$\mathbf{N = 2}$}\\
%          \multirow{2}{*}{2} & GPT-4 & 78.2 & 100.0 & 0.0 & 2.000 \\
%          & Cascaded Selective Evaluation & 85.1 & 64.7 & 90.5 & 1.288 \\
%          \midrule
%          % \rowcolor{gray!25}\multicolumn{5}{c}{$\mathbf{N = 3}$}\\
%          \multirow{2}{*}{3} & GPT-4 & 78.1 & 100.0 & 0.0 & 3.000 \\
%          & Cascaded Selective Evaluation & 85.3 & 67.9 & 90.3 & 1.920 \\
%          \midrule
%          % \rowcolor{gray!25}\multicolumn{5}{c}{$\mathbf{N = 5}$}\\
%          \multirow{2}{*}{5} & GPT-4 & 78.5 & 100.0 & 0.0 & 5.000 \\
%          & Cascaded Selective Evaluation & 85.2 & 70.2 & 90.6 & 3.131 \\
%          \bottomrule
%     \end{tabular}
% }
% \vspace{-5pt}
% \label{tab:impact_of_number_of_simulated_annotators}
% \end{table*}
Table \ref{tab:analyzing_abstention_policy} presents the results. The average inter-annotator agreement is 0.815 ($\sigma^2 = 0.031$) for abstained samples and 0.902 ($\sigma^2 = 0.025$) for evaluated samples, a statistically significant difference in two-sample t-test with $p < 1e-8$. This is in contrast with both the length ratio and token overlap, for which the differences between the two sets are not significant ($p > 0.10$). In fact, for token overlap, the abstained examples exhibit higher ROUGE-L on average than the evaluated samples. Overall, these results show that the instances abstained by LLM judges tend to be more subjective even for humans (with no evidence of reliance on some spurious heuristics), indicating that the confidence elicited by Simulated Annotators closely aligns with that of human annotators.

\subsection{Evaluation under Distribution Shift} Our test procedure in \S \ref{sec:providing_human_agreement_guarantee} assumes that the calibration set $D_\textit{cal}$ is sampled i.i.d. from $P(x, y_\textit{human})$. In real-world scenarios this may not be the case, because we often only have access to generations from a set of known models for calibration, while in the test time, we need to evaluate outputs from unknown models. In \S\ref{app:evaluation_under_distribution_shift}, we empirically analyze whether our method provides risk control even under this distribution shift. First, we randomly divide ChatArena into two disjoint sets such that there is no overlap between the evaluated models in each set. Then, we induce distribution shift by sampling $D_\textit{cal}$ from one set and testing instances in another set. We follow the same setup as \S \ref{sec:evaluating_chat_assistants}, and run experiments for 1000 random splits. As shown in Table \ref{tab:evaluation_under_distribution_shift}, despite a small degradation in coverage, our method guarantees high human agreement for more than 90\% of the time, consistently across all tested levels of $\alpha$. The result demonstrates that \method maintains its reliability even under the realistic distribution shift.



\subsection{Impact of Number of Simulated Annotators}\label{sec:impact_of_number_of_simulated_annotators} 
Next, we analyze the impact of number of simulated annotators for \method. In Table \ref{tab:impact_of_number_of_simulated_annotators}, we report the results on ChatArena using Simulated Annotators as confidence measure with $N = 5, 3, 2, 1$. We compare the result against using GPT-4 with the same prompt, but without abstention. Along with the guarantee success rate and coverage, we report relative API cost for calling OpenAI models, where the cost for full evaluation with GPT-4 is set to 1. The results suggest that (1) using larger number of simulated annotators leads to consistently better coverage, but (2) even with a small number of simulated annotators, our method can still achieve high human agreement while reducing the evaluation cost by up to 40\% compared to GPT-4 without abstention.

\subsection{Impact of Judge Model Composition} We study whether \method can be done entirely without GPT-4. We use a substantially weaker cascades with \textit{Mistral-7B-instruct-v0.2}, \textit{Mixtral-8$\times$7b-instruct} \citep{mixtral}, and GPT-3.5 as judge models (\textit{weaker} cascades). We compare the result against (1) zero-shot GPT-4 without abstention, and (2) \method using the original cascades, with $N = 1$ for better cost (\textit{stronger} cascades). We set the target agreement level $1 - \alpha = 0.8$, a higher level than what is achievable by GPT-4 without abstention (Figure \ref{fig:chatarena_results}). 

The results in Table \ref{tab:impact_of_judge_composition} reveal an interesting finding: even the weaker cascades of judge models ensure a satisfactory level of human agreement, by balancing the the trade-off between the evaluation cost and coverage, instead of compromising the precision. Unlike conventional LLM-based evaluation, where one has to sacrifice accuracy by using a weaker judge, our method allows practitioners to consistently achieve their target level of human agreement. Depending on the requirements, one can opt for stronger cascades for better coverage or weaker cascades for lower costs, all while maintaining this guarantee. Additionally, both configurations of \method significantly reduce evaluation costs compared to using GPT-4, achieving savings of up to 78.5\% with stronger cascades and 87.4\% with weaker cascades.

% \begin{table*}[t]\centering
% \vspace{-15pt}
% \caption{Caption TBU}
% \resizebox{.93\textwidth}{!}{
%     \begin{tabular}{lcccc}\toprule
%          \multirow{2}{*}{\textbf{Method}} & \multirow{2}{*}{\shortstack{\textbf{Empirical Human}\\\textbf{Agreement (\%)}}} & \multirow{2}{*}{\textbf{Coverage (\%)}} & \multirow{2}{*}{\shortstack{\textbf{Guarantee Success}\\\textbf{Rate (\%)}}} & \multirow{2}{*}{\shortstack{\textbf{Relative }\\\textbf{API Cost}}} \\
%          & & & & \\
%          \midrule
%          GPT-4 & 77.8 & 100.0 & 13.9 & 1.000 \\
%          \midrule
%          \multirow{2}{*}{\shortstack[l]{Cascaded Selective Evaluation \\ (\textit{stronger cascades)}}} & \multirow{2}{*}{80.2} & \multirow{2}{*}{88.4} & \multirow{2}{*}{90.2} & \multirow{2}{*}{0.215} \\
%          & & & & \\
%          \midrule
%          \multirow{2}{*}{\shortstack[l]{Cascaded Selective Evaluation \\ (\textit{weaker cascades)}}} & \multirow{2}{*}{80.3} & \multirow{2}{*}{74.4} & \multirow{2}{*}{90.1} & \multirow{2}{*}{0.126} \\
%          & & & & \\
%          \bottomrule
%     \end{tabular}
% }
% \vspace{-5pt}
% \label{tab:impact_of_judge_composition}
% \end{table*}

\begin{table*}[t]\centering
\vspace{-15pt}
\caption{Impact of judge model composition on ChatArena, with $1 - \alpha = 0.8$. \textit{Weaker} cascades use \textit{Mistral-7B}, \textit{Mixtral-8$\times$7B}, and \textit{GPT-3.5} as judge models. Stronger cascades use \textit{GPT-4} instead of \textit{Mixtral-8$\times$7B}. \textbf{Our method guarantees human agreement even with the weaker cascades, while only using 12.6\% of the evaluation cost for GPT-4.}}
\vspace{-5pt}
\resizebox{.93\textwidth}{!}{
    \begin{tabular}{lcccc}\toprule
         \multirow{2}{*}{\textbf{Method}} & \multirow{2}{*}{\shortstack{\textbf{Empirical Human}\\\textbf{Agreement (\%)}}} & \multirow{2}{*}{\textbf{Coverage (\%)}} & \multirow{2}{*}{\shortstack{\textbf{Guarantee Success}\\\textbf{Rate (\%)}}} & \multirow{2}{*}{\shortstack{\textbf{Relative }\\\textbf{API Cost}}} \\
         & & & & \\
         \midrule
         GPT-4 & 77.8 & 100.0 & 13.9 & 1.000 \\
         Cascaded Selective Evaluation \textit{(stronger)} & 80.2 & 77.6 & 90.5 & 0.215 \\
         % & & & & \\
         Cascaded Selective Evaluation \textit{(weaker)} & 80.3 & 68.3 & 90.8 & 0.126 \\
         % & & & & \\
         \bottomrule
    \end{tabular}
}
\vspace{-10pt}
\label{tab:impact_of_judge_composition}
\end{table*}